\subsection{A fixed Hilbert series with unbounded canonical parameter degrees}

\begin{proposition}[P-F22: the Gorenstein condition does not repair the transfer]
\label{prop:gorensteinperiodcollapse}
Let $q_0=1$, $q_1=2$, and $q_{j+1}=3q_j-q_{j-1}$ for $j\ge1$.
For $j\ge1$ put $b=q_j$, $c=q_{j+1}$ and
\[
 Q_{b,c}=\{(x,y)\in\mathbb R^2:x\ge-1,\ y\ge-1,
                       b^2x+c^2y\le1\}.
\]
Its Ehrhart ring
\[
 A_{b,c}=\mathbb C[\{(x,y,s)\in\mathbb Z^3:
                         s\ge0,\ (x,y)\in sQ_{b,c}\}]
\]
is a positively graded normal Cohen--Macaulay Gorenstein domain
of Krull dimension three with rational singularities. For every
pair in this sequence it has the same Hilbert function and series,
\[
 \dim(A_{b,c})_s=\binom{3s+2}{2}
       =\frac92s^2+\frac92s+1\quad(s\ge0),
 \qquad
 \operatorname{Hilb}_{A_{b,c}}(t)=\frac{1+7t+t^2}{(1-t)^3}.
 \tag{P-F22a}
\]
Its positivity and saturation indices are zero. Nevertheless its
lexicographically minimal nondecreasing homogeneous parameter-degree
tuple is exactly $(1,b,c)$. In particular its canonical maximal
parameter degree is $c$, which tends to infinity with $j$.
No member of this family is identified here with an ordinary
irreducible-inner plethysm multiplicity algebra, and factoriality
is not asserted.
\end{proposition}

\begin{proof}
The recurrence preserves the identity
\[
 1+b^2+c^2=3bc.
 \tag{P-F22b}
\]
Indeed it holds for $(b,c)=(1,2)$, and substitution of
$(c,3c-b)$ into the left minus right side leaves
$1+b^2+c^2-3bc$. Consecutive terms are relatively prime because
$\gcd(c,3c-b)=\gcd(c,b)$, and neither is divisible by three
because reduction of the recurrence modulo three sends
$(b,c)$ to $(c,-b)$. For our pairs $b\ge2$ and
\[
 2b<c<3b.
 \tag{P-F22c}
\]
This holds for $(2,5)$; from $0<b<c$ the next ratio is
$(3c-b)/c=3-b/c$, and induction preserves the displayed bounds.
In particular the sequence is unbounded.

We first count the smaller rational triangle
\[
 P_{b,c}=\operatorname{conv}\{(0,0),(c/b,0),(0,b/c)\}.
\]
We use the following precise published period-collapse criterion.
For $u=q/p$, $v=s/r$ in lowest terms,
Cristofaro-Gardiner--Li--Stanley, \emph{New examples of period
collapse}, Corollary~1.5, printed p.~4 of arXiv
\texttt{1509.01887v1}, states that the triangle with intercepts
$1/u,1/v$ has polynomial Ehrhart function if
\[
 s\mid p,\quad p\mid(rq+1),\quad
 \gcd((rq+1)/p,s)=1,
\]
and
\[
 q\mid r,\quad r\mid(sp+1),\quad
 \gcd((sp+1)/r,q)=1.
\]
This is the source's explicit rational-triangle criterion, not
its separate theorem concerning irrational triangles. Its primary
text is \url{https://arxiv.org/pdf/1509.01887}.

Here take $(p,q,r,s)=(c,b,b,c)$. The divisibilities follow from
\[
 b^2+1=c(3b-c),\qquad c^2+1=b(3c-b).
\]
The two gcd conditions are
$\gcd(3b-c,c)=1$ and $\gcd(3c-b,b)=1$, which follow from
$\gcd(b,c)=\gcd(3,bc)=1$. All hypotheses of the criterion are
therefore satisfied. The Ehrhart function of $P_{b,c}$ is a
quadratic polynomial with leading coefficient its area $1/2$
and constant coefficient one. By (P-F22c), its only lattice
points at dilation one are $(0,0),(1,0),(2,0)$. Its value there
is three, so this polynomial is
\[
 L_{P_{b,c}}(s)=\tfrac12s^2+\tfrac32s+1
                    =\binom{s+2}{2}.
 \tag{P-F22d}
\]
In obtaining the value at zero we use the usual Ehrhart
quasi-polynomial, which counts the single point $0P_{b,c}$;
the criterion says its constituents coincide.

Equation (P-F22b) gives
\[
 Q_{b,c}=3P_{b,c}-(1,1).
\]
Translation by $(-s,-s)$ is integral at every integer dilation,
so (P-F22d) yields $L_{Q_{b,c}}(s)=L_{P_{b,c}}(3s)$.
This proves the Hilbert function in (P-F22a). Summing its
quadratic polynomial, or multiplying its first three differences
by $(1-t)^3$, gives the displayed Hilbert series. Its ordinary
polynomial coefficients are positive and all nonnegative
integer values are positive, so both indices are zero.

The primitive rays of the cone over $Q_{b,c}$ are
\[
 v_0=(-1,-1,1),\quad
 v_1=(3c-b,-b,b),\quad
 v_2=(-c,3b-c,c).
 \tag{P-F22e}
\]
Primitivity of the last two follows from
$\gcd(3c,b)=\gcd(3b,c)=1$. The cone is simplicial and
full-dimensional. As in the full argument of P-F21, subtracting
integer parts of its nonnegative ray coordinates gives a unique
decomposition into a lattice point in the half-open fundamental
parallelepiped plus a nonnegative integral combination of the
three rays. Thus $A_{b,c}$ is finite free over
$\mathbb C[X^{v_0},X^{v_1},X^{v_2}]$, whose generator degrees
are $(1,b,c)$. This proves finite generation, Krull dimension
three, Cohen--Macaulayness, and existence of a parameter sequence
with those degrees. Normality follows by expressing the ring as
the intersection of the three integrally closed halfspace
semigroup rings for
\[
 \ell_1(x,y,s)=x+s,\quad
 \ell_2(x,y,s)=y+s,\quad
 \ell_3(x,y,s)=s-b^2x-c^2y.
 \tag{P-F22f}
\]
These are primitive integral facet forms. Their nonnegative
intersection is precisely the cone: its inequalities imply
$s\ge0$, since $x,y\ge-s$ and
$b^2x+c^2y\le s$ give $(1+b^2+c^2)s\ge0$; at $s=0$
they force $x=y=0$.

The Gorenstein assertion has a direct lattice check. With
$e=(0,0,1)$, each $\ell_i(e)=1$. Hence an integral point is
strictly inside the cone exactly when its translate by $-e$
is an integral point of the cone:
\[
 \operatorname{int}(C_{b,c})\cap\mathbb Z^3
       =e+(C_{b,c}\cap\mathbb Z^3).
 \tag{P-F22g}
\]
Danilov, \emph{The geometry of toric varieties}, Section~4.6,
printed pp.~111--112, identifies the canonical module of a
normal affine toric ring with the ideal spanned by its interior
monomials. Equation (P-F22g) identifies that module with the
principal ideal $X^eA_{b,c}$. Since the ring is Cohen--Macaulay,
it is Gorenstein. Its rational singularities follow from the
toric resolution and vanishing in the same source, Section~8.2
and Proposition~8.5.1, printed pp.~124--125. Both exact locators
are available in the primary scan
\url{https://www.math.purdue.edu/~arapura/danilov.pdf}.

We finally prove the canonical degree assertion; mere existence
of $(1,b,c)$ would not by itself be a lexicographic comparison.
Let $F$ be the two-dimensional cone face spanned by $v_1,v_2$,
equivalently the face $\ell_3=0$. Its degree-$s$ lattice points
correspond, after $u=x+s$, $v=y+s$, to nonnegative integers
satisfying
\[
 b^2u+c^2v=3bc\,s.
\]
Reduction modulo $b$ gives $b\mid v$, and reduction modulo $c$
gives $c\mid u$. Writing $u=cr$, $v=bt$ therefore reduces the
condition exactly to
\[
 br+ct=3s,\qquad r,t\in\mathbb Z_{\ge0}.
 \tag{P-F22h}
\]
There is no solution for $0<s<b$. If $t\ge2$, then
$br+ct\ge2c>4b>3s$. If $t=1$ and $r\ge1$, it is at least
$b+c>3b>3s$. If $t=1,r=0$, then $c=3s$, contrary to
$3\nmid c$. If $t=0$, then $br=3s$, and $\gcd(b,3)=1$
forces $b\mid s$, again impossible.

The monomial face quotient $A_{b,c}\twoheadrightarrow
\mathbb C[F\cap\mathbb Z^3]$ is a graded domain of dimension
two. Every positive-degree element of degree less than $b$ maps
to zero in it. A parameter sequence cannot have two such
elements: its remaining single element would cut a
two-dimensional finitely generated domain to a ring of
dimension at least one, or leave it unchanged, contradicting
finite-dimensionality of the quotient by all three parameters.
Here a nonzero homogeneous positive-degree element is not a
unit, and the dimension assertion is the principal ideal theorem.
Consequently the second entry of every nondecreasing
parameter-degree tuple is at least $b$.

The extreme-ray quotient at $v_2$ is $\mathbb C[z]$ with
$\deg z=c$. If a parameter tuple begins $(1,b)$, both these
parameters map to zero there because $1,b<c$. The third must
map nontrivially, so its degree is a positive multiple of $c$.
The first entry is always at least one. Together with the
existing tuple $(1,b,c)$, these facts prove that this tuple is
lexicographically minimal and its maximal degree is exactly $c$.
\end{proof}

\paragraph{What this additional example excludes.}
The construction uses a standard Markov and Fibonacci period-collapse
family and the published rational-triangle criterion; no historical
priority is claimed. Unlike P-F21, here the entire Hilbert series,
including every numerator coefficient, is independent of the
unbounded parameter degree. Thus even knowing the complete
Hilbert function, fixed Krull dimension, normality,
Cohen--Macaulayness, Gorensteinness, and rational singularities
does not bound the canonical modular index of a general graded
domain. All examples already have a positive form of denominator
weights one. This identifies a further invalid general route
from P-F19 to the weaker positive-form clause; it supplies no
ordinary-plethysm counterexample to that weaker clause.
